\documentclass[a4paper,11pt]{article}
\usepackage[utf8]{inputenc}
\usepackage[T1]{fontenc}
\usepackage[english]{babel}
\usepackage{geometry}
\geometry{left=2cm,right=2cm,top=2cm,bottom=2cm}

\usepackage{lmodern}
\usepackage{xcolor}
\definecolor{titleblue}{RGB}{0,51,140}
\usepackage{hyperref}
\usepackage{titlesec}

\titleformat{\section}{\color{titleblue}\large\bfseries}{\thesection}{1em}{}

\begin{document}

\begin{center}
    {\LARGE\bfseries\color{titleblue} Research Statement}\\[6pt]
    {\large Structural Methods for Mixed Model/Data Digital Twin Engineering}\\[4pt]
    \textbf{Ismail AISSAOUI} $|$ State Engineer in AI
\end{center}

\bigskip

\section{Introduction and Motivation}
Modern engineering digital twins are built upon large-scale physical models derived from first principles. However, the complexity of these models often leads to significant discrepancies when compared to real-world data, originating from parameter uncertainty, sensor drift, or structural modeling deficiencies. This research statement proposes a framework to address these "reality gaps" using structural analysis of differential-algebraic equation (DAE) systems, enabling the precise localization and correction of model/data mismatches.

\section{Current Research: Physics-Informed Digital Twins}
In my current work at Sonatrach, I have developed a generative surrogate model for gas turbine monitoring using \textbf{Mamba-SSM} and \textbf{xLSTM} architectures. My contribution lies in the design of \textbf{Physics-Informed Neural Network (PINN)} loss functions that enforce thermodynamic consistency. This work has highlighted a critical challenge: when a discrepancy occurs, is it due to the data, the parameters, or the fundamental physical equations? My experience has demonstrated the need for a formal structural diagnostic layer that precedes or accompanies the data assimilation process.

\section{Proposed Research Directions}
Building upon the objectives of the \textbf{Structural Methods} project at Inria Rennes, I propose to investigate the following axes:

\subsection{1. Graph-Based Structural Analysis for DAEs}
I intend to leverage my background in \textbf{Graph Neural Networks (GNNs)} to develop scalable algorithms for the structural analysis of large-scale DAE systems. By representing equations and variables as bipartite graphs, I will explore the computation of Minimal Structurally Overdetermined (MSO) subsystems. These subsystems act as "parity spaces," providing a mathematical foundation for detecting inconsistencies without requiring the full reachability analysis of the system.

\subsection{2. Localizing Model/Data Mismatches}
The core of my research will focus on using the residuals generated by solved MSO subsystems to localize deficiencies. By performing statistical analysis on these residuals, I aim to distinguish between sensor errors (data-driven) and structural deficiencies in the underlying Modelica or port-Hamiltonian models. This targeted diagnosis is essential for the efficient maintenance of complex industrial twins.

\subsection{3. Scalability for Large Sparse Systems}
Real-world industrial systems can contain thousands of equations. I propose to optimize the structural analysis workflow for large sparse systems by leveraging high-performance graph algorithms. The goal is to ensure that the diagnostic tools remain efficient for real-time or near-real-time synchronization within the digital twin lifecycle.

\section{Alignment with the EDT Program}
My background in relational modeling and industrial AI implementation aligns with the interdisciplinary goals of the **EDT program**. I am eager to contribute to the **Artemis platform** by providing a robust structural diagnostic layer that ensures the "truthfulness" of the hybrid models used across the consortium.

\section{Conclusion}
The goal of my PhD will be to transform digital twin engineering from ad-hoc data assimilation to a rigorous, structurally-aware diagnostic process. By bridging the gap between graph theory, DAE structural analysis, and hybrid modeling, I aim to provide the tools necessary for the next generation of reliable and self-diagnosing digital twins.

\end{document}
